% Prover/verifier protocol box (house style).  Toy protocol: prove that a
% committed polynomial f vanishes on a domain H via its quotient q = f / Z_H.  Columns: prover computation on the left,
% verifier on the right, messages as labelled horizontal arrows in the middle,
% rounds top-to-bottom, verifier checks numbered so a proof can cite them.
% For a paper body, cryptocode's \pseudocode gives the same layout
% (skills/paper-playbook/sections/08_main_construction.md).
\documentclass[border=3pt]{standalone}
\usepackage[T1]{fontenc}
\usepackage{lmodern}   % LNCS; for ACM/IEEE use \usepackage{newtxtext,newtxmath}
\usepackage{amsmath,amssymb}
\usepackage{crypto-tikz}
\newcommand{\alg}[1]{\mathsf{#1}}
\begin{document}
\begin{tikzpicture}[x=1cm,y=-0.62cm]
  \node[party] at (0,0) {Prover $\mathcal{P}(x; w)$};
  \node[party] at (8,0) {Verifier $\mathcal{V}(x)$};
  \draw[crgrid, line width=0.4pt] (-1.6,0.55) -- (9.6,0.55);
  % round 1: commitment
  \node[anchor=west] at (-1.6,1.4) {$f \gets \alg{Encode}(x, w),\ q \gets f / Z_H$};
  \node[anchor=west] at (-1.6,2.3) {$C_f \gets \alg{Com}(f),\ C_q \gets \alg{Com}(q)$};
  \draw[msg] (2.2,3.5) -- node[above, font=\footnotesize] {$C_f,\ C_q$} (5.8,3.5);
  % round 2: challenge
  \node[anchor=east] at (9.6,4.4) {$\rho \xleftarrow{\$} \mathbb{K}$};
  \draw[msg] (5.8,5.3) -- node[above, font=\footnotesize] {$\rho$} (2.2,5.3);
  % round 3: evaluations and opening proof
  \node[anchor=west] at (-1.6,6.2) {$y_f \gets f(\rho),\ y_q \gets q(\rho)$};
  \node[anchor=west] at (-1.6,7.1) {$\pi \gets \alg{Open}(f, q; \rho)$};
  \draw[msg] (2.2,8.3) -- node[above, font=\footnotesize] {$y_f,\ y_q,\ \pi$} (5.8,8.3);
  % checks, numbered so the proof can cite them
  \node[anchor=east] at (9.6,9.2) {(1)\ $\alg{Vrfy}(C_f, C_q; \rho, y_f, y_q, \pi) = 1$};
  \node[anchor=east] at (9.6,10.1) {(2)\ $y_f = Z_H(\rho)\, y_q$};
  \node[anchor=east, text=crink2, font=\footnotesize] at (9.6,11.0) {accept iff (1) and (2)};
  % round markers
  \foreach \y/\r in {3.5/1, 5.3/2, 8.3/3}
    \node[op, anchor=west] at (-1.6,\y) {round \r};
  % frame
  \draw[crink2, line width=0.4pt, rounded corners=1.5pt] (-1.8,-0.7) rectangle (9.8,11.6);
\end{tikzpicture}
\end{document}
